\bookchapter{Kapitel 9: Mobiles Arbeiten und Lernen}

\chapterintro{Früher habe ich auch Arbeiten am Smartphone nur als „Notbehelf“ betrachtet – bis ich auf einer Dienstreise den Laptop vergaß und Angebotsänderung wie Vertragsunterzeichnung komplett am Handy erledigte. Ehrlich gesagt liegt der Engpass mobiler Arbeit selten am Gerät selbst, sondern daran, dass man seine Fähigkeiten und Grenzen nicht kennt.}

\booksection{Dokumente bearbeiten: Schreiben und Gestalten in der Handfläche}

Dokumente auf dem Smartphone-Bildschirm zu bearbeiten fühlt sich wirklich anders an als am Computer. Die größten Probleme: kleiner Bildschirm, dicke Finger, schnell verrutschtes Layout. Mit ein paar Tricks lässt sich das Erlebnis aber deutlich verbessern:

\begin{itemize}
  \item \textbf{WPS Office}: Am besten an den chinesischen Markt angepasst, vollständig kompatibel mit Microsoft-Formaten (docx/xlsx/pptx), mit reichhaltigem integriertem PDF-Werkzeugkasten. Besonders praktisch finde ich die Funktion „Sprachnotiz“ – wenn du beim Lesen eines Vertrags nicht tippen kannst, sprich einfach ein paar Anmerkungen ins Handy, und sie werden automatisch als Textkommentar in das Dokument eingefügt. Außerdem bringt WPS riesige Mengen kostenloser Vorlagen mit – beim Lebenslauf oder beim Präsentations-PPT sparst du so viel Zeit.
  \item \textbf{Microsoft 365}: Wenn deine Firma eine lizenzierte Office-Version nutzt, ist das die beste Wahl. Die Formatkompatibilität mit der Desktopversion ist praktisch perfekt – das leichte Verrutschen komplexer Layouts, das WPS gelegentlich zeigt, tritt hier nicht auf. Dank tiefer OneDrive-Integration ist ein Dokument, das du gerade am PC bearbeitest, auf dem Smartphone schon in der aktuellen Version.
  \item \textbf{Shimo Docs / Tencent Docs}: Diese beiden sind die Musterbeispiele für „Online-Kollaboration“. Sie eignen sich, wenn mehrere Personen gleichzeitig ein Dokument bearbeiten und jede Änderung in Echtzeit sichtbar wird. Ich plane mit Freunden Reisepläne in Shimo Docs – wenn zwei gleichzeitig die Tabelle ausfüllen, ist das deutlich effizienter als allein.
\end{itemize}

\textbf{Praktischer Workflow}: Meine jetzige Gewohnheit: wichtige Dateien in OneDrive, alltägliches Bearbeiten mit WPS, Teamkollaboration in Shimo Docs. Seit meine Dateien in der Cloud liegen, gab es die Peinlichkeit „USB-Stick vergessen“ kein einziges Mal mehr.

\booksection{Notiz-Apps: vom spontanen Notieren zum Wissenssystem}

Beim Notizen machen wollen viele sofort ein hochkomplexes System aufbauen und halten dann keine drei Tage durch. Mein Rat: erst leicht, dann schwer – entwickle zuerst die Gewohnheit des Aufschreibens und kümmere dich danach um die Verwaltung.

\begin{itemize}
  \item \textbf{Spontanes Notieren – die Boardmittel reichen}: Die iOS-Notizen-App startet extrem schnell und unterstützt Dokumentenscan, Handschrift und Sprachaufnahme. Auch die Notiz-Apps der Android-Hersteller sind gut und lassen sich schon vom Sperrbildschirm aus schnell öffnen. Diese Werkzeuge eignen sich für spontane Ideen, Einkaufslisten und Besprechungspunkte. Entscheidend ist der Schnellstart – eingefallen, notiert, ohne auf das Laden der App zu warten.
  \item \textbf{Später lesen – Informationen sammeln}: Die „Schwebende Fenster"-Funktion von WeChat ist eigentlich ein verstecktes Meisterwerk. Wenn Sie einen guten Artikel sehen, aber keine Zeit haben, ihn zu lesen, tippen Sie oben rechts auf „Schwebendes Fenster", und er wird zu einem kleinen Punkt am Bildschirmrand minimiert, den Sie jederzeit wieder öffnen können. Das ist viel bequemer, als ihn an den „Dateiübertragungsassistenten" weiterzuleiten. Außerdem können Tools wie Cubox oder Pocket gute Inhalte plattformübergreifend sammeln, die Sie irgendwo gesehen haben.
  \item \textbf{System aufbauen – Notion und Feishu}: Wenn Sie das Gefühl haben, dass verstreute Notizen schwer zu finden sind, ist es Zeit für Notion. Dieses Tool ist wirklich flexibel – Sie können eine Datenbank für Lese-notizen anlegen, die automatisch nach Buchtitel, Autor und Bewertung sortiert wird; oder ein „Projekt-Board" erstellen, um den Fortschritt von Arbeitsaufgaben zu visualisieren. Die „Mehrdimensionalen Tabellen" von Feishu sind ähnlich und im Inland schneller erreichbar. Ich verwalte mit Notion alle meine Schreibmaterialien, Lernnotizen und Projektfortschritte – seit über drei Jahren.
\end{itemize}

\textbf{Eine praktische Notizstrategie}: Meine eigene Gewohnheit ist das „Drei-Teile-Set“ – spontane Einfälle in der System-Notiz (Sekunden), zu lesende Artikel im WeChat-Schwebefenster (bequem), langfristig gesammeltes Wissen in Notion (systematisch). Versuchen Sie nicht, alles sofort in ein einziges Tool zu stopfen – das führt eher dazu, dass man aufgibt.

\booksection{Scannen und OCR: Die Papierwelt ins Smartphone holen}

Das Scannen mit dem Smartphone kann, richtig eingesetzt, viele Gänge zum Copyshop ersparen. Als ich das letzte Mal ein Visum beantragte, brauchte ich elektronische Versionen einer ganzen Reihe von Unterlagen – alles mit dem Smartphone erledigt.

\begin{itemize}
  \item \textbf{iPhone-Nutzer}: Die in iOS-Notizen integrierte Scanfunktion ist bereits sehr gut. Öffnen Sie Notizen tippen Sie auf das Kamera-Symbol wählen Sie „Dokument scannen". Richten Sie das Dokument aus, es wird automatisch erkannt, der Rahmen erfasst und die Perspektive korrigiert. Nach der Aufnahme können Sie die Farbe anpassen – der Modus „Schwarz-Weiß" liefert Ergebnisse, die einem Scangerät am nächsten kommen, „Graustufen" eignet sich für Materialien mit Abbildungen.
  \item \textbf{Android-Nutzer}: Viele Markenkameras haben einen „Dokument-Scan"-Modus. Falls nicht, laden Sie Microsoft Lens aus dem App-Store herunter (kostenlos und mit OneDrive verknüpft).
  \item \textbf{OCR-Texterkennung}: Wenn Sie aus einem gescannten Dokument den Text zur Bearbeitung extrahieren möchten, benötigen Sie OCR. Die in WeChat integrierte Funktion „Text extrahieren" ist am einfachsten – langes Drücken auf das Bild genügt, und die Erkennungsgenauigkeit ist ordentlich. Für größere Mengen eignen sich die App Baimiao oder Quanneng Scan King. Für die chinesische Zeichenerkennung liefert Baidus OCR-Engine gute Ergebnisse; die Bild-zu-Text-Funktion von WPS ist ebenfalls sehr genau.
  \item \textbf{Scan-Tipps}: Fotografieren Sie an einem Ort mit ausreichendem Licht, am besten mit einer Tischlampe seitlich, um Schatten der Hände auf dem Dokument zu vermeiden. Wenn das Dokument geknickt ist, glätten Sie es erst mit der Hand. Mehrseitige Dokumente können nacheinander gescannt und zu einem einzigen PDF zusammengefügt werden, ohne jede Seite einzeln zu speichern.
\end{itemize}

\booksection{Cloud-Speicher wählen: Es gibt nicht den Besten, nur den für Sie Passenden}

Bei der Wahl des Cloud-Speichers habe ich lange überlegt und schließlich erkannt: Es kommt darauf an, was Ihr digitales Hauptleben ausmacht.

\begin{itemize}
  \item \textbf{Apple-Ökosystem iCloud}: Die nahtlose Synchronisierung lässt Sie seine Existenz kaum spüren. Ein Foto, das Sie auf dem Smartphone aufnehmen, ist Minuten später automatisch auf dem Computer. Und Apple nimmt Datenschutz am strengsten. Aber die Nachteile sind offensichtlich: Das Teilen mit Android- oder Windows-Nutzern ist umständlich, kostenlos gibt es nur 5 GB, Erweiterungen kosten Geld.
  \item \textbf{Windows-Nutzer OneDrive}: Wenn Sie einen Windows-Computer + Office 365 nutzen, ist OneDrive die natürliche Wahl. Speichern Sie ein Dokument in Word, landet es automatisch in der Cloud. 1 TB Speicherplatz in Kombination mit dem Office-Abo ist sehr wirtschaftlich. Die automatische Foto-Backup-Funktion der Smartphone-App ist ebenfalls nützlich.
  \item \textbf{Teilen im Inland Baidu Netdisk}: Ganz ehrlich, die Drosselung der Download-Geschwindigkeit von Baidu Netdisk wird oft kritisiert, aber sein Sharing-Ökosystem ist unübertroffen. Eine große Menge an Lernmaterialien und Medien im chinesischsprachigen Internet wird über Baidu Netdisk geteilt. Wer nur gelegentlich Dateien von anderen herunterlädt, kommt auch ohne Mitgliedschaft aus. Wer jedoch viele Ressourcen sammelt, sollte die Mitgliedschaft abschließen – die gesparte Zeit ist das Geld wert.
  \item \textbf{Alibaba Cloud Drive / Quark Netdisk}: Blogger mit Smartphone-Fotografie-Tutorials folgen.
\end{itemize}

\textbf{Meine Kombination}: iCloud für persönliche Fotos und Alltagsdateien (Datenschutz zuerst), Baidu Netdisk für Ressourcen und empfangene Freigaben (Ökosystem zuerst), OneDrive für Arbeitsdokumente (Office-Integration zuerst). Drei Cloud-Laufwerke, jedes mit seiner Aufgabe, ohne sich gegenseitig zu stören.

\booksection{Ein realistisches Szenario mobiles Arbeiten}

Erzählen Sie eine eigene wahre Geschichte. Einmal auf einer Geschäftsreise in einer anderen Stadt, der Kunde wollte kurzfristig Vertragskonditionen ändern. Mein Workflow sah so aus:

\begin{enumerate}
  \item Im Hotel erhielt ich die Änderungswünsche des Kunden per WeChat (Textnachricht und ein Foto mit handschriftlichen Anmerkungen).
  \item Ich öffnete OneDrive, fand das Vertragsoriginal, bearbeitete es mit WPS und änderte die Klauseln entsprechend den WeChat-Nachrichten des Kunden Punkt für Punkt.
  \item Nach dem Ändern exportierte ich mit der WPS-Funktion „Als PDF ausgeben" und schickte sie zur Bestätigung in die WeChat-Gruppe des Kunden.
  \item Der Kunde hatte eine zweite Runde Anmerkungen. Ich nutzte das Werkzeug „Stift" in WPS, um Änderungsstellen direkt im PDF zu markieren, und ergänzte Sprachnotizen, um die Änderungsideen zu erläutern.
  \item Nach Bestätigung der Fehlerlosigkeit nutzte ich die WPS-Funktion „Elektronische Signatur" zur Online-Unterzeichnung und schickte sie dem Gegenüber zur gleichen Unterzeichnung.
  \item Zum Schluss speicherte ich das unterzeichnete Vertrags-PDF jeweils in Baidu Netdisk und iCloud und notierte gleichzeitig einen Eintrag im Arbeitsjournal in Notion.
\end{enumerate}

Der gesamte Prozess dauerte vom Anfang bis zum Ende etwa vierzig Minuten, ganz ohne Computer. Ganz ehrlich: Vor drei Jahren hätte ich nicht geglaubt, dass das Smartphone diese Dinge erledigen kann, aber heute geht es tatsächlich. Der Schlüssel ist, diese Werkzeuge im Voraus kennenzulernen, anstatt erst im Bedarfsfall hektisch Downloads und Registrierungen durchzuführen.

\begin{tcbpractice}{Probieren Sie es aus: Bauen Sie Ihren mobilen Arbeitsworkflow auf}
\begin{enumerate}
  \item Wählen Sie Ihre Haupt-Dokumenten-App (WPS oder Microsoft 365), melden Sie sich an und verknüpfen Sie den Cloud-Speicher.
  \item Öffnen Sie Notizen oder Notizzettel, scannen Sie ein Dokument, exportieren Sie es als PDF und senden Sie es an sich selbst.
  \item Wählen Sie eine Cloud-Speicher-App und richten Sie das automatische Backup von Smartphone-Fotos ein (am besten „Nur über WLAN sichern", um Datenvolumen zu sparen).
  \item Erstellen Sie in Notion oder Feishu eine neue Seite und versuchen Sie, eine einfache „Arbeitsjournal"-Vorlage zu erstellen.
\end{enumerate}

\end{tcbpractice}
